\section*{अध्याय \ref{ch:g7:oxygen-in-the-water} --- पानी में ऑक्सीजन}

\begin{solution}{exo:g7:oxygen-in-the-water:1}
प्रति लीटर लगभग \qty{10}{mg} --- यानी एक लीटर हवा जितनी ढोती है, उससे
मोटे तौर पर तीस गुना कम।
\end{solution}

\begin{solution}{exo:g7:oxygen-in-the-water:2}
ऊपर की हवा, जो सतह पर घुलती है --- सबसे तेज़ वहाँ, जहाँ पानी हिलता,
छपकता और झाग बनाता है; और नीचे के \omterm{def:g5:food-webs:roles}{उत्पादक}, जो दिन में \omterm{def:g7:what-is-respiration:respiration}{ऑक्सीजन} छोड़ते
हैं --- सबसे तेज़ तेज़ \omterm{def:g6:exploring-our-environment:conditions}{प्रकाश} में, भरपूर हरी बढ़वार के साथ।
\end{solution}

\begin{solution}{exo:g7:oxygen-in-the-water:3}
\omterm{def:g1:animals-around-us:animal}{जंतुओं} का \omterm{def:g7:what-is-respiration:respiration}{श्वसन}, रात में ख़ुद \omterm{def:g1:plants-around-us:plant}{पौधों} का \omterm{def:g7:what-is-respiration:respiration}{श्वसन}, और अपघटकों का --- जो मरे
हुए पदार्थ की भरमार होने पर बहुत बड़ा हो जाता है।
\end{solution}

\begin{solution}{exo:g7:oxygen-in-the-water:4}
गरम पानी \omterm{def:g7:oxygen-in-the-water:dissolved}{घुली हुई ऑक्सीजन} कम थामता है (छत गिर जाती है), जबकि हर साँस
लेने वाले का इंजन तेज़ चलता है (माँग चढ़ जाती है): एक ही साथ आपूर्ति कम
और ज़रूरत ज़्यादा।
\end{solution}

\begin{solution}{exo:g7:oxygen-in-the-water:5}
दिन में \omterm{def:g5:food-webs:roles}{उत्पादकों} का निर्माण \omterm{def:g7:what-is-respiration:respiration}{ऑक्सीजन} उससे तेज़ जोड़ता है जितनी \omterm{def:g7:what-is-respiration:respiration}{श्वसन}
घटाता है --- वक्र चढ़कर दोपहर बाद के शिखर पर पहुँचता है। रात भर हर साँस
लेने वाला भंडार पर हाथ डालता है और निर्माण बिलकुल नहीं होता --- वक्र
फिसलकर सूर्योदय से ठीक पहले अपने सबसे नीचे आ जाता है।
\end{solution}

\begin{solution}{exo:g7:oxygen-in-the-water:6}
भरपूर: पत्थरों के नीचे ट्राउट (या ट्राउट के बच्चे) और मेफ़्लाई के डिंभ।
ग़रीब: सिर्फ़ सहनशील \omterm{def:g5:biodiversity:species}{प्रजातियाँ} --- कार्प, तालाब के घोंघे, और सतह से
हवा लाने वाले नलीदार डिंभ।
\end{solution}

\begin{solution}{exo:g7:oxygen-in-the-water:7}
रात का \omterm{def:g7:what-is-respiration:respiration}{श्वसन} --- \omterm{prop:g3:sorting-living-things:groups}{मछलियाँ}, \omterm{def:g1:plants-around-us:plant}{पौधे} और कीचड़ के \omterm{def:g6:decomposers-and-soil:decomposers}{अपघटक}, सब मिलकर --- तालाब को
सूर्योदय से ठीक पहले उसके न्यूनतम तक घटा देता है, इसलिए बोझ से लदे तालाब
अपनी \omterm{prop:g3:sorting-living-things:groups}{मछलियों} का दम भोर में ही घोंटते हैं। इलाज: पहले आपूर्तिकर्ता से,
सतह को हिलाओ --- रात भर फ़व्वारा या झरना चलाओ; दूसरे आपूर्तिकर्ता से,
घास-फूस और कीचड़ पतला करो ताकि रात का \omterm{def:g7:what-is-respiration:respiration}{श्वसन} घटे (और दिन का उत्पादन
सिर्फ़ सतह की चटाई तक नहीं, पानी तक पहुँचे)।
\end{solution}

\begin{solution}{exo:g7:oxygen-in-the-water:8}
कचरा भीतर --- \omterm{def:g6:decomposers-and-soil:decomposers}{अपघटक} दावत उड़ाते और बढ़ते हैं --- उनका जुटा हुआ \omterm{def:g7:what-is-respiration:respiration}{श्वसन}
\omterm{def:g7:what-is-respiration:respiration}{ऑक्सीजन} सोख लेता है --- माँग करने वाली \omterm{def:g5:biodiversity:species}{प्रजातियाँ} ग़ायब और सहनशील
क़ाबिज़। नीचे की ओर कचरा आख़िर खप जाता है, पथरीली उथल-पुथल पानी को फिर
हवा दे देती है, और माँग करने वाली \omterm{def:g5:biodiversity:species}{प्रजातियाँ} लौट आती हैं: घटाने वाला बंद
हुआ, आपूर्तिकर्ता आगे निकल आए।
\end{solution}

\begin{solution}{exo:g7:oxygen-in-the-water:9}
मापक यंत्र एक पल दर्ज करता है; बाशिंदे इतिहास दर्ज करते हैं। कोई माँग
करने वाली \omterm{def:g5:biodiversity:species}{प्रजाति} किसी ठिकाने पर तभी जीती है जब पानी उसके पूरे बसेरे भर
अच्छा बना रहा हो --- उन भोरों और उन बुरे हफ़्तों समेत, जिन्हें कोई
अकेला चक्कर नहीं पकड़ता। गिनती पूरे मौसम का न्यूनतम है, जो \omterm{def:g1:animals-around-us:animal}{जंतुओं} में
लिखा हुआ है।
\end{solution}

\begin{solution}{exo:g7:oxygen-in-the-water:10}
मटर-सूप जैसी हरियाली हरी अकेली \omterm{def:g5:first-look-at-cells:cell}{कोशिकाओं} की एक विशाल आबादी है: दिन में
वे उत्पादन करती हैं, पर हर गरम रात पूरा सूप साँस लेता है --- और वह गरम
ठहरा हफ़्ता छत गिरा देता है तथा सतह की भराई रोक देता है। भोर वाला
न्यूनतम शून्य की ओर बढ़ता है: यानी मछली-मौत का मौसम। फ़व्वारा रात और
छोटे घंटों में सबसे ज़ोर से चलाओ --- दिन ख़ुद अपना ख़्याल रख लेगा।
\end{solution}

\begin{solution}{exo:g7:oxygen-in-the-water:11}
ट्राउट माँग करने वाली \omterm{def:g5:biodiversity:species}{प्रजाति} है: ठंडी नदियाँ उसे वह ऊँची छत देती हैं
जो उसे चाहिए; और छोटे बाँध तथा छपकती नालियाँ भीड़ भरी \omterm{prop:g3:sorting-living-things:groups}{मछलियों} के घटाने
के ख़िलाफ़ भराई लगातार बनाए रखते हैं। गरम ठहरी झील छत और भराई, दोनों पर
एक साथ नाकाम है --- वह कार्प का इलाक़ा है, ट्राउट का नहीं।
\end{solution}

\begin{solution}{exo:g7:oxygen-in-the-water:12}
मुद्रा: पानी के नीचे हर साँस लेने वाला उसी एक घुले हुए भंडार पर हाथ
डालता है --- कहाँ कौन फलता-फूलता है, इसकी क़ीमत \omterm{def:g7:what-is-respiration:respiration}{ऑक्सीजन} में चुकती है
(\cref{prop:g7:oxygen-in-the-water:map})। मुद्रा की आपूर्ति: पानी कितना
थाम सकता है यह तापमान तय करता है --- ठंड छत ऊपर उठाती है, गरमी नीचे
गिराती है (\cref{prop:g7:oxygen-in-the-water:drains})। सबसे बड़े
ख़र्चीले: मरा हुआ पदार्थ मिल जाए तो अपघटकों का जुटा हुआ \omterm{def:g7:what-is-respiration:respiration}{श्वसन} सबसे
ज़्यादा घटाता है --- यही डेयरी के पाइप का सबक़ है
(\cref{ex:g7:oxygen-in-the-water:waste})।
\end{solution}

\begin{solution}{pb:g7:oxygen-in-the-water:1}
\textbf{1.} क1 ठंडा है --- ऊँची छत --- और बाँध से मथा हुआ --- सतह की
भराई पूरी रफ़्तार पर: इसलिए छत के पास की पढ़तें। क2 गरम है --- गिरी हुई
छत --- और बिना हिला हुआ: उसके आपूर्तिकर्ता बस एक मामूली पढ़त ही थाम
सकते हैं।

\textbf{2.} क1 की माँग करने वाली \omterm{def:g5:biodiversity:species}{प्रजातियाँ} भरपूर और टिकाऊ \omterm{def:g7:what-is-respiration:respiration}{ऑक्सीजन} की
गवाही देती हैं; क2 का सहनशील जमघट --- कार्प, घोंघे, नली वाले --- ऐसे
पानी की निशानी है जो कम से कम अपने बुरे घंटों में ग़रीब पड़ जाता है।

\textbf{3.} अपघटकों का \omterm{def:g7:what-is-respiration:respiration}{श्वसन}: डेयरी का कार्बनिक-भरपूर बहाव उनकी एक जुटी
हुई आबादी को खिलाता है, और उनकी दावत भंडार को तापमान अकेले जितना समझा
सकता, उससे कहीं नीचे घटा देती है।

\textbf{4.} दो किलोमीटर की खपत ने कचरा चुका दिया है --- दावत ख़त्म और
ख़र्चीले पतले पड़ गए --- जबकि पथरीली उथल-पुथल ने \omterm{def:g7:what-is-respiration:respiration}{ऑक्सीजन} फिर से मथकर
भीतर पहुँचा दी है: दोनों आपूर्तिकर्ता बहाल, घटाने वाला बंद, और गिनती इस
सुधार की गवाही देती है।

\textbf{5.} एक रोज़ की लहर: भोर का गड्ढा \qty{4}{mg/L} पर, दोपहर बाद की
चोटी \qty{9}{mg/L} पर। जगह बदलने वाले सौदागर: दिन का निर्माण (दिन में
जोड़ता हुआ) और चौबीसों घंटे चलता \omterm{def:g7:what-is-respiration:respiration}{श्वसन} (रात भर अकेला घटाता हुआ)।

\textbf{6.} सर्दी की ठंड छत ऊपर उठा देती है, इसलिए पानी ज़्यादा थामता
है; और ठंड हर साँस लेने वाले को --- \omterm{prop:g3:sorting-living-things:groups}{मछलियाँ}, \omterm{def:g1:plants-around-us:plant}{पौधे}, कीचड़ के \omterm{def:g6:decomposers-and-soil:decomposers}{अपघटक} ---
धीमा कर देती है, इसलिए घटाने वाले सिकुड़ जाते हैं। और ज़्यादा टिकाऊ भी,
क्योंकि \omterm{def:g5:food-webs:roles}{उत्पादकों} की रोज़ की लहर छोटी रहती है।

\textbf{7.} पहले क3 (पहले से घटा हुआ, और गरमी उसके अपघटकों को तेज़ करती
है); दूसरे क2 (नीची छत, ठहरा पानी, भोर के गड्ढे); तीसरे क4 (सुधर रहा है,
पर मुसीबत के नीचे की ओर); और आख़िर में क1 (ठंडा और मथा हुआ --- बाँध मौसम
चाहे जो हो, भराई जारी रखता है)।

\textbf{8.} सूर्योदय से पहले का समय रोज़ का न्यूनतम है: तब की पढ़तें वह
सबसे बुरा हाल दर्ज करती हैं जिससे नदी के जीवन को गुज़रना ही पड़ता है।
दोपहर बाद वाली समय-सारणी तो साल को उसके सबसे अच्छे पलों के नाम दर्ज कर
देती।

\textbf{9.} मापक यंत्र कुछ ही हफ़्तों में सुधर जाता है --- अपघटकों का
भोजन रुकते ही घटाने वाला बंद हो जाता है। गिनती ज़्यादा धीरे और क्रम से
सुधरती है: पहले सहनशील जमघट पतला पड़ता है, फिर बहकर और उड़कर आने वाले
बसने वाले (\cref{ch:g6:how-plants-spread} का तर्क, \omterm{def:g1:animals-around-us:animal}{जंतुओं} वाला रूप)
माँग करने वाली \omterm{def:g5:biodiversity:species}{प्रजातियों} को फिर से बो देते हैं --- पहले ऊपर की ओर से
मेफ़्लाई के डिंभ, और सबसे बाद में ट्राउट।

\textbf{10.} वह छाया उस हिस्से को ठंडा भी रखती है --- और ठंड ही तो छत
है। काटने से दिन का उत्पादन बढ़ेगा, पर छत गिरेगी और गरम हुए पानी में हर
घटाने वाला तेज़ होगा; धीमे हिस्से पर नुक़सान शायद फ़ायदे से भारी है।
जवाब: \omterm{def:g1:plants-around-us:tree}{पेड़} रहने दो, पानी को हिलाओ।

\textbf{11.} गिनती पर: मेफ़्लाई के डिंभ अच्छे मौसम का नाटक नहीं कर सकते
--- वे उसकी हर भोर जीकर आए हैं। दोपहर बाद वाली मामूली पढ़त एक ही पल है,
शायद कोई ढीला दिन; डिंभ पूरे मौसम का हलफ़िया लेखा हैं।

\textbf{12.} ``नदी को ठंडा, मथा हुआ और बिना बोझ का रखो, फिर \omterm{def:g7:what-is-respiration:respiration}{ऑक्सीजन} ---
और उसे साँस में लेने वाला सब कुछ --- अपना ख़्याल ख़ुद रख लेगा।''
\end{solution}
